% Round 18. Higher-center branch. Candidate; independent review required.
\section{Full Green coverage in every even dimension for the hyperplane models}
\label{r18ev:section}

The hyperplane construction is not confined to the fourfold middle
degree.  In even generic dimension its double strata have numerical
groups of constant rank in all interior degrees.  Multiplication by
the pullback of the original hyperplane class identifies their
displayed bases.  A single sufficiently large ample perturbation
then gives all the injectivities required by DG18's recursive
criterion, including degrees above the middle.  No primitive
Hodge positivity on a component is used.

Assume the complete native hypotheses of HG18's
dimension-independent ordered hyperplane model:
\[
       X=V\!\left(\prod_{i=1}^d\ell_i+\pi Q\right)
                                      \subset\mathbf P_K^{n+1},
 \qquad n\ge4,\quad d\ge3,
 \tag{EV18.1}\label{r18ev:model}
\]
the reductions of the hyperplanes are in geometric linear general
position, $\bar Q$ cuts every positive-dimensional linear stratum
smoothly and avoids the zero-dimensional strata, and the generic
hypersurface is smooth and geometrically integral.  Let
$\mathcal X/R$ be the actual native reduced strict semistable model
constructed there.  In this section assume
\[
                              n=2k\quad\hbox{is even}.
 \tag{EV18.2}\label{r18ev:parity}
\]
Put $m=n-1$.  All Chow and numerical groups have rational
coefficients; every field and degree is native.
Let $f:\mathcal X\to\mathbf P_R^{n+1}$ be the model morphism and
write $h=f^*c_1(\mathcal O(1))$, using the same letter for its
restrictions.  This class is globally generated; it is not
assumed ample on the resolved strata.

\subsection{All numerical degrees on every double}

For an edge $E_{ij}$, $i<j$, HG18 identifies it as the iterated
codimension-two blowup of $\mathbf P_k^m$ in the descending
sequence of centers
\[
              Z_{ij,l}=H_{ijl}\cap V(\bar Q),\qquad l>j.
 \tag{EV18.3}\label{r18ev:centers}
\]
Every center, including its strict transform at the relevant step,
is isomorphic to the indicated smooth hypersurface in
$\mathbf P_k^{n-2}$, of odd dimension $n-3=m-2$.
Let $f_{ij,l}$ be the exceptional divisor class of that step,
pulled back through all later steps, as in HG18.

\begin{lemma}[Rank-one comparison on the centers, and comparison on the edges]
\label{r18ev:comparison}
For every center in \eqref{r18ev:centers} and every
$0\le v\le m-2$,
\[
                   N^v(Z_{ij,l})_{\mathbf Q}
                         =\mathbf Q(h|_{Z_{ij,l}})^v,
 \tag{EV18.4}\label{r18ev:center-rank}
\]
and every actual native numerically zero rational cycle on that
center has zero geometric $\ell$-adic class, for every
$\ell\ne\operatorname{char}k$.
Every double $E_{ij}$ has the same numerical-to-geometric
comparison in every codimension.  Its numerical groups have
the explicit bases
\[
 \begin{aligned}
 N^0(E_{ij})&=\mathbf Q\,1,\\
 N^u(E_{ij})&=
       \mathbf Q h_{ij}^u
       \ \oplus\!
       \bigoplus_{l>j}\mathbf Q\,f_{ij,l}h_{ij}^{u-1}
                                  &&(1\le u\le m-1),\\
 N^m(E_{ij})&=\mathbf Q h_{ij}^m.
 \end{aligned}
 \tag{EV18.5}\label{r18ev:edge-bases}
\]
All these assertions remain valid after every finite extension
of the residue field, with the new cycles on that field.
\end{lemma}

\begin{proof}
The center is a smooth positive-dimensional hypersurface and
is geometrically integral by the argument already checked in
HG18.  Its dimension is odd.  We apply only the
odd-dimensional complete-intersection comparison proposition
of frozen CI17, not CI17's full rank-one-boundary model
criterion.  That proposition follows from geometric weak
Lefschetz and Poincare duality: every even cohomology group
of an odd-dimensional smooth complete intersection is the
line generated by the corresponding hyperplane power.
For an actual native cycle $z$ of codimension $v$, the
coefficient of that line is the rational number
\[
      \frac{\deg_k(z h^{m-2-v})}{\deg_k(h^{m-2})}.
 \tag{EV18.6}\label{r18ev:rational-coefficient}
\]
Thus numerical zero implies geometric homological zero, and
every numerical class is the stated hyperplane multiple.
The argument applies afresh after every finite residue extension,
since it uses the full geometric even-cohomology line rather
than a possibly changing Frobenius-fixed subspace.

For a codimension-two smooth blowup $q:M'\to M$ with
exceptional embedding $j:J\to M'$ and projective-line map
$p:J\to Z$, the actual Chow decomposition is
\[
       \operatorname{CH}^u(M')
           =q^*\operatorname{CH}^u(M)
                    \oplus j_*p^*\operatorname{CH}^{u-1}(Z).
 \tag{EV18.7}\label{r18ev:blowup}
\]
HC18 proves its numerical decomposition and coefficient
extraction by proper pushforward, restriction and exceptional
self-intersection.  In particular numerical zero of an
actual class in \eqref{r18ev:blowup} implies numerical zero
of both actual coefficients.
If the old variety and the center satisfy comparison, those
coefficients have zero geometric cycle class.  Functoriality
of geometric cycle classes under smooth-variety pullback
and Gysin push then gives zero class for the original cycle.
Thus all-degree comparison is preserved by the displayed
blowup, without asserting a cycle-class map on a numerical
group before this has been proved.

Start with projective space and apply this argument successively.
The additional numerical summand is $N^{u-1}(Z)$, which has
rank one exactly for $1\le u\le m-1$ and is zero in the two
endpoint degrees.  The pullback of the ambient hyperplane
class to each unchanged center is its original hyperplane
class.  The added generator is therefore precisely
$f_{ij,l}h_{ij}^{u-1}$, by the actual projection formula.
This proves the bases and comparison on the final edge.
All coefficient operations and decompositions are native
and remain valid after the stated scalar extensions.
\end{proof}

\begin{lemma}[A common ample class which is an isomorphism in all interior edge degrees]
\label{r18ev:ample-isomorphisms}
There exists one relatively ample integral divisor class $H$
on $\mathcal X$ such that
\[
            H:W^u\xrightarrow{\ \sim\ }W^{u+1}
                         \qquad(1\le u\le n-3).
 \tag{EV18.8}\label{r18ev:edge-isomorphisms}
\]
\end{lemma}

\begin{proof}
By the displayed actual generators in
\eqref{r18ev:edge-bases}, multiplication by the ambient class
$h$ sends each basis vector in degree $u$ to its identically
indexed basis vector in degree $u+1$, for
$1\le u\le m-2=n-3$:
\[
           h\cdot h^u=h^{u+1},\qquad
           h\cdot(f_{ij,l}h^{u-1})=f_{ij,l}h^u.
 \tag{EV18.9}\label{r18ev:ambient-isomorphism}
\]
Both target classes are in the basis in this range.
The formula uses the pulled-back exceptional generators,
so it is unchanged by later blowups, including at old
higher-stratum intersections.

Choose any relatively ample integral line bundle
$\mathcal L$ on the projective model and set
\[
                      H_M=c_1(\mathcal L)+M h,\qquad M\ge0.
\]
Because $h$ is globally generated, every such integral
$H_M$ is relatively ample.
On each edge and in each indicated degree, its multiplication
matrix in \eqref{r18ev:edge-bases} is
$M I+C_{ij,u}$ for a fixed rational square matrix $C_{ij,u}$.
Its determinant is a monic polynomial in $M$.
There are finitely many edges and finitely many indicated
degrees.  Choose one sufficiently large nonnegative integer
outside the finite union of roots of these polynomials.
Then every matrix is invertible and gives
\eqref{r18ev:edge-isomorphisms} with $H=H_M$.
No numerical hard Lefschetz or middle Hodge conjecture has
been used; the proof is the explicit basis calculation.
\end{proof}

\subsection{Recursive local exactness and full ordinary actual G}

\begin{theorem}[All required middle defects vanish in even dimension]
\label{r18ev:defects}
For the model in \eqref{r18ev:model}--\eqref{r18ev:parity},
\[
                         U_r=0\qquad(1\le r\le n-2).
 \tag{EV18.10}\label{r18ev:Urzero}
\]
For every relevant degree $a$ one also has
$K_{\mathrm{comp},\ell}^a=\ker\rho_a$.
Consequently every specified geometrically homologically
trivial rational cycle on $X$ has strong ordinary-actual G
in every codimension, on the complement of its original
precise support and against all compact ordinary actual tests.
\end{theorem}

\begin{proof}
DG18 gives $U_1=0$ unconditionally.  HG18 gives $U_2=0$
on the actual hyperplane model, retaining all its triple
and higher strata.  These vanishings are statements about
the fixed numerical maps, independent of any subsequent
choice of an ample class for the recursive argument.

Choose the one $H$ of Lemma~\ref{r18ev:ample-isomorphisms}.
For $3\le r\le n-2$, the two multiplications
\[
    H:W^{r-2}\xrightarrow{\ \sim\ }W^{r-1},
       \qquad
    H:W^{r-1}\xrightarrow{\ \sim\ }W^r
 \tag{EV18.11}\label{r18ev:induction-pair}
\]
lie in the range \eqref{r18ev:edge-isomorphisms}.
Assume inductively that $U_{r-1}=0$.
DG18's recursive injection is applicable by the second
isomorphism, and its target is zero by the first:
\[
       U_r\hookrightarrow
       \frac{W^{r-1}}{\operatorname{im}\rho_{r-1}
                      +H W^{r-2}+\ker\sigma_r}=0.
\]
This proves \eqref{r18ev:Urzero}.  When $n=4$ the
induction range is empty and the conclusion is exactly
the established $U_1,U_2$ result.

For the representative gate, take any numerically compatible
tuple and choose actual rational Chow representatives on
its components.  Their differences on the doubles are
numerically zero, and Lemma~\ref{r18ev:comparison} makes
their geometric cycle classes zero.  They therefore give
compatible actual representatives.  The converse is the
frozen homological-implies-numerical assertion, so
$K_{\mathrm{comp},\ell}^a=\ker\rho_a$.
This is the precise CL17 comparison mechanism.

For $2\le r\le n-2$, set $a=n-r$.
The two gates of DQ17 now vanish, so ER17 gives zero
obstruction for every specified geometrically homologically
trivial codimension-$r$ cycle.
The frozen exact-support reconstruction gives its Green form
with the full original actual test quantifier.
The remaining positive codimensions are $r=1,n-1,n$;
these are the already frozen divisor, one-cycle and
zero-cycle cases.  Hence the local G domain is the complete
geometrically homologically trivial rational Chow group in
every codimension.
\end{proof}

\paragraph{What is, and is not, asserted about model quotients.}
The proof establishes the numerical defects in
\eqref{r18ev:Urzero}, and hence the corresponding refined
quotient vanishings.  It does not infer vanishing of every
model quotient merely from full G in the endpoint cycle
degrees: the sufficient model criterion is not necessary.
In particular the multiplication from numerical curves on
an edge to its top degree need not be injective; it is
outside \eqref{r18ev:edge-isomorphisms} and is never used.

\begin{corollary}[Deep local families in every even dimension]
\label{r18ev:deep}
If $d\ge n+2$, the even-dimensional family above has full
local G and the extreme monodromy map of HM18 has rank
$\binom{d-1}{n+1}>0$.  Thus these are not merely new
models of potentially good-reduction generic fibers.
For every even $n\ge4$, EX18's explicit equation
\[
                 \prod_{i=0}^{n+1}X_i
                       +\pi\sum_{i=0}^{n+1}X_i^{n+2}=0,
 \qquad p\nmid n+2,
 \tag{EV18.12}\label{r18ev:explicit}
\]
has full local G in all codimensions and rank-one extreme
monodromy.  These Green assertions reapply to arbitrary
new cycles after every unramified finite extension of $K$.
\end{corollary}

\begin{proof}
Combine the theorem with the separately checked model and
monodromy hypotheses of HM18 and EX18.
Under an unramified extension, $\pi$ remains a uniformizer
and all geometric transversality conditions persist.
Lemma~\ref{r18ev:comparison} and the ample-matrix construction
apply over the new native residue field, so the theorem
applies anew to every cycle over that field.
No descent of those new cycles to the original field is
assumed.
\end{proof}

\subsection{The same full conclusion for nonlinear factors}

\begin{corollary}[All even-dimensional nonlinear-factor models have full G]
\label{r18ev:nonlinear}
Assume all the hypotheses of NF18:
\[
 X=V(f_1\cdots f_s+\pi Q)\subset\mathbf P_K^{n+1},
 \qquad s\ge3,\quad \deg f_i=d_i\ge1,\quad
 \deg Q=\sum_i d_i.
 \tag{EV18.13}\label{r18ev:nonlinear-model}
\]
Thus the reduced factor hypersurfaces form a geometric
simple-normal-crossings divisor, every nonempty intersection
is a smooth complete intersection of the expected dimension,
$\bar Q$ cuts every positive-dimensional such intersection
transversely and avoids all zero-dimensional intersections,
and $X$ is smooth and geometrically integral.
If $n\ge4$ is even, the actual ordered model of NF18 satisfies
$U_r=0$ for $1\le r\le n-2$ and
$K_{\mathrm{comp},\ell}^a=\ker\rho_a$ in every degree.
Every specified geometrically homologically trivial rational
cycle on $X$ therefore satisfies strong ordinary-actual G
in every codimension.  The assertion applies afresh to
all cycles after every finite unramified extension of $K$.
\end{corollary}

\begin{proof}
NF18 supplies the actual regular reduced strict semistable
model and identifies its edge $E_{ij}$ as the iterated
codimension-two blowup of
$H_{ij}=H_i\cap H_j$ along the unchanged strict transforms
of $Z_{ij,l}=H_i\cap H_j\cap H_l\cap V(\bar Q)$,
for $l>j$.
Now $H_{ij}$ is a smooth complete intersection of odd
dimension $n-1$, and each $Z_{ij,l}$ is a smooth complete
intersection of odd dimension $n-3$.  Both are geometrically
integral by the NF18 hypotheses and connectedness argument.
Frozen CI17's odd-dimensional comparison proposition gives,
in every numerical degree, the one-dimensional hyperplane
subspace and numerical-to-geometric comparison on both
the starting variety and each center.  The rational
coefficient is computed by intersecting with the complementary
hyperplane power and dividing by the nonzero native
complete-intersection degree, exactly as in
\eqref{r18ev:rational-coefficient}.

The actual blowup decomposition
\eqref{r18ev:blowup} therefore gives
\eqref{r18ev:edge-bases}, with the exceptional classes of
NF18 in place of $f_{ij,l}$, and gives numerical-to-geometric
comparison in every degree on every final edge.
The varying positive factor degrees affect the nonzero
top intersection numbers; they do not change these basis
vectors or the exact multiplication
$h(h^u)=h^{u+1}$,
$h(e_{ij,l}h^{u-1})=e_{ij,l}h^u$
in the interior degree range.
Consequently the single ample perturbation argument
$H_M=c_1(\mathcal L)+Mh$ has the same square matrices
$MI+C_{ij,u}$ and proves
\eqref{r18ev:edge-isomorphisms} on this nonlinear model.

NF18 supplies $U_2=0$ using its actual triple table, whose
exceptional block is still $-I$; DG18 supplies $U_1=0$.
The identical induction \eqref{r18ev:induction-pair} gives
every $U_r$ asserted here.
The all-degree edge comparison closes the CL17
representative gate, and DQ17, ER17 and frozen reconstruction
give the middle-cycle Green forms.  The frozen endpoint
cycle theorems complete all codimensions.
All hypotheses, complete-intersection comparisons and
coefficient decompositions reapply after each unramified
extension, so the conclusion is for its full new cycle
groups and does not assume descent of their cycles.
\end{proof}

\paragraph{Dependencies, literature and remaining scope.}
The geometry is HG18's native stratum description, including
its identification of all successive centers.
The only complete-intersection comparison invoked is the
frozen odd-dimensional proposition of CI17.  Its primary
cohomological input is weak Lefschetz and Poincare duality;
CI17 records Deligne, \emph{La conjecture de Weil. II},
Section 4.1.6, as the checked source for weak Lefschetz.
The blowup decomposition and numerical coefficient extraction
are HC18, with the standard projective-bundle and excess
formulas.  The recursive injection is DG18, and the
representative and ordinary-actual conclusions use the
frozen CL17, DQ17, ER17 and specified-cycle reconstruction.

No component middle-positivity assertion and no arbitrary
proper actual pushforward enters the proof.
No complete-intersection rank-one model criterion is applied
to the blown-up edges: their extra exceptional directions
are essential and permit genuine higher strata.
An arbitrary ramified extension changes the model problem
and is not covered merely by the unramified statement.
Likewise a number-field height assertion still requires
placewise full G at every finite place; this local theorem
alone does not verify those other places.
There is no claim of general G for arbitrary varieties or
of historical priority for the classical degeneration.
